\section{Motivation}

Music choice is contextual. A track that suits exercise may be unwelcome during
a quiet break. People use music to regulate mood and activation
~\cite{juslin2008,schaefer2013}. Most recommender systems prioritize durable
signals such as listening history, favorite artists, similar users, and item
features. Those signals characterize usual preference and capture momentary
fit only indirectly, a known challenge in contextual music recommendation
~\cite{kaminskas2012,schedl2018}. A listener may like a song in general without
wanting it next.

Explicit mood selection can address that gap, but it also interrupts an
otherwise continuous experience. Automatic adaptation offers another route.
Its observations, however, are uncertain. Facial behavior depends on pose,
lighting, culture, context, and deliberate expression
~\cite{barrett2019,aviezer2012,kuppens2017}. Cardiac timing varies with
respiration, posture, movement, fitness, medication, and sensor quality
~\cite{shaffer2017,laborde2017,quigley2024}. Neither channel directly identifies
emotion.

\emph{Vibe Shuffle} uses these uncertain observations only as within-listener
changes that may guide the next song. During setup, the listener
provides short personal camera and cardiac references. Later observations are
expressed as changes from those references. The horizontal axis represents
relative Valence evidence; the vertical axis represents relative Arousal
evidence. Each point falls into a Tense, Energetic, Melancholic, or Calm
recommendation region. These names identify catalog regions, not the
listener's emotional state.

The catalog uses song Valence and Energy coordinates. Vibe Shuffle averages
the recent trajectory, retains unplayed tracks in its region, and selects the
nearest eligible track. Random draws from the same remaining catalog without
the trajectory. The adaptive chain is explicit: personal references, relative
observations, trajectory, region, and nearest available song.

The study compares this policy with Random selection from the same fixed
catalog. Its primary outcome is participant-rated mood fit, analyzed through
the paired difference and the uncertainty around it. Liking is a secondary
measure of general song preference. The comparison rests on three reproducible
elements:
\begin{itemize}
  \item Personal reference periods center camera and cardiac observations
  before the relative Valence and Arousal coordinates are calculated.
  \item Event-driven observations form a trajectory whose mean identifies a
  catalog region and whose distance ranking yields the nearest unplayed track.
  \item The exploratory crossover estimates mood-fit and liking differences
  under matched catalog, exposure, and rating procedures.
\end{itemize}

Mood-fit estimates favored Vibe Shuffle, whereas liking barely changed; both
intervals crossed zero. A service-level evaluation would require validated
sensor mappings, complete trial records, preregistration, and a dynamic catalog.

\begin{figure*}[t]
  \centering
  \includegraphics[width=\textwidth]{fig1_affect_space.pdf}
  \caption{Panels (a)--(b) are original redrawings adapted from Helmholz et
  al.~\cite{helmholz2019}, who
  operationalized Russell's circumplex~\cite{russell1980} and Thayer's
  mood--arousal model~\cite{thayer1991}. Panel (c) is an original schematic
  adapted from the group-level frequency distinction reported by Balaji et
  al.~\cite{balaji2025}; it contains no empirical observations.}
  \Description{Three square conceptual panels show a continuous
  valence-arousal plane, a categorical affect vocabulary, and a
  coherence-frequency schematic.}
  \label{fig:affect-space}
\end{figure*}
